\section{Conclusions}\label{sec:conclusions}

The fact that the gauge field generated by the gradient flow
does not require renormalization is a consequence of the
locality, the symmetries and the parabolic nature
of the flow equation.
In particular, in the
associated field theory in $D+1$ dimensions,
bulk counterterms are excluded simply
because the evolution of the bulk fields is retarded at large times
and thus described by tree diagrams.
The absence of
boundary counterterms other than those needed for the renormalization of
the theory at flow time zero is however
non-trivial and can only be shown
using power-counting and the BRS symmetry.

In presence of matter fields, the situation is essentially
unchanged as long as only the gauge field is evolved in time.
The finiteness of the correlation functions of the field at positive
flow time is therefore still guaranteed.
Including all or some of the matter fields in the flow is however
an interesting option that remains to be explored.
